%:strip
% marginnotes.tex: Provide marginnotes support
% Copyright (c) 2021 by SIL International 
% written by Martin Hosken
% 
% This optional plugin (see ptx-plugins) extends the basic diglot engine to
% provide marginal notes.
%
% Permission is hereby granted, free of charge, to any person obtaining
% a copy of this software and associated documentation files (the  
% "Software"), to deal in the Software without restriction, including  
% without limitation the rights to use, copy, modify, merge, publish,  
% distribute, sublicense, and/or sell copies of the Software, and to  
% permit persons to whom the Software is furnished to do so, subject to  
% the following conditions:
%
% The above copyright notice and this permission notice shall be  
% included in all copies or substantial portions of the Software.
%
% THE SOFTWARE IS PROVIDED "AS IS", WITHOUT WARRANTY OF ANY KIND,  
% EXPRESS OR IMPLIED, INCLUDING BUT NOT LIMITED TO THE WARRANTIES OF  
% MERCHANTABILITY, FITNESS FOR A PARTICULAR PURPOSE AND  
% NONINFRINGEMENT. IN NO EVENT SHALL SIL INTERNATIONAL BE LIABLE FOR  
% ANY CLAIM, DAMAGES OR OTHER LIABILITY, WHETHER IN AN ACTION OF  
% CONTRACT, TORT OR OTHERWISE, ARISING FROM, OUT OF OR IN CONNECTION  
% WITH THE SOFTWARE OR THE USE OR OTHER DEALINGS IN THE SOFTWARE.
%
% Except as contained in this notice, the name of SIL International  
% shall not be used in advertising or otherwise to promote the sale,  
% use or other dealings in this Software without prior written  
% authorization from SIL International.
%%%%%%%%%%%%%%%%%%%%%%%%%%%%%%%%%%%%%%%%%%%%%%%%%%%%%%%%%%%%%%%%%%%%%%%

% TeX does not allow these to be executed inside a conditional:
\newif\ifm@rginnotesopen 
\newif\ifm@rginalnotesw@rn
\newwrite\m@rginnotes
%\plugin@startif{marginnotes}

\def\@marginnote#1#2#3#4#5#6#7#8{%
% #1 = reference, #2 = marker, #3 = side, #4 = valign, #5 = distance from text,
% #6 = width, #7 = height, #8 = depth, #9 = x-shift, #10 = y-shift
  \@m@rginnote{#1}{#2}% chain to collect parameters over #9
}

\def\@m@rginnote#1#2#3#4#5#6#7{%
% #1 = reference, #2 = marker, #3 = x-shift, #4 = y-shift, #5 = pageno, #6 = xpos, #7 = ypos
  \x@\xdef\csname marginnote-#1@#2:xshift\endcsname{#3}%
  \x@\xdef\csname marginnote-#1@#2:yshift\endcsname{#4}%
  \x@\xdef\csname marginnote-#1@#2:pageno\endcsname{#5}%
  % identify which column it actually was/is in
  \dimen0=#6 sp \adjustpdfXsavepos \dimen1=\dimexpr \textwidth + \SideMarginFactor \MarginUnit\relax\divide \dimen1 2
  \ifBindingGutter\ifodd #5\else\advance\dimen1 \BindingGutter\fi\fi
  \ifdim\dimen0 > \dimen1 \ifRTL\edef\c@l{0}\else\edef\c@l{1}\fi\else\ifRTL\edef\c@l{1}\else\edef\c@l{0}\fi\fi
  \x@\xdef\csname marginnote-#1@#2:colno\endcsname{\c@l}%
  \trace{M}{defining marginnote-#1@#2:pageno = #5, xpos=\the\dimen0, col=\c@l, based on \the\dimen0 >\the\dimen1 (smf=\SideMarginFactor x\the\MarginUnit, tw=\the\textwidth, bg=\the\BindingGutter)}%
}

% Pushed notes (see docs/dev/marginnotes_overflow.md). A note with a push line
% is held in \m@rginpending and drawn at the top of its target column, one stack
% per side, at the first output for its target page.
\def\@marginnotepush#1#2#3#4#5{%
% #1 = reference, #2 = marker, #3 = target page, #4 = target column, #5 = side
  \x@\xdef\csname marginnote-#1@#2:push\endcsname{{#3}{#4}{#5}}%
  % write the push line back, so Python can see it has been carried out
  \xdef\m@rginpushes{\m@rginpushes\write\m@rginnotes{\string\@marginnotepush{#1}{#2}{#3}{#4}{#5}}}%
}
\def\m@rginpushes{}
\newif\ifm@rginheld
\newcount\m@rginserial
\newbox\m@rginpending       % held notes: \penalty<serial> then the lapped note box
\newbox\m@rginnewpending
\x@\newbox\csname m@rginstack@0@left\endcsname \x@\newbox\csname m@rginstack@0@right\endcsname
\x@\newbox\csname m@rginstack@1@left\endcsname \x@\newbox\csname m@rginstack@1@right\endcsname
\def\m@rginstack#1{\csname m@rginstack@#1\endcsname}
\def\m@rginreleasedpage{}

\def\m@rginhold{% box0 = lapped note
  \global\advance\m@rginserial by 1
  \global\setbox\m@rginpending=\vbox{\unvbox\m@rginpending\penalty\m@rginserial\nointerlineskip\box0}%
  \x@\xdef\csname mnheld@\the\m@rginserial\endcsname{{\m@rginkey}\csname marginnote-\m@rginkey:push\endcsname}%
  \x@\xdef\csname mnlatest@\m@rginkey\endcsname{\the\m@rginserial}% a paragraph set twice holds its notes twice
}

% At the first output for a page, move held notes due by this page into the
% stacks. Walking from the bottom and prepending keeps the original order.
\def\m@rginovfout{%
  \ifnum\pageno=0\m@rginreleasedpage\else
    \xdef\m@rginreleasedpage{\the\pageno}%
    \setbox0=\vbox{\unvbox\m@rginpending\m@rginw@lk}%
    \global\setbox\m@rginpending=\box\m@rginnewpending
  \fi
}
\def\m@rginw@lk{%
  \setbox0=\lastbox
  \ifvoid0 \else
    \count255=\lastpenalty \unpenalty
    \x@\x@\x@\m@rginrele@se\csname mnheld@\the\count255\endcsname
    \x@\m@rginw@lk
  \fi
}
\def\m@rginrele@se#1#2#3#4{% #1 = key, #2 = page, #3 = column, #4 = side; box0 = note, \count255 = serial
  \x@\ifnum\csname mnlatest@#1\endcsname=\count255
    \ifnum#2>\pageno
      \global\setbox\m@rginnewpending=\vbox{\penalty\count255\nointerlineskip\box0\unvbox\m@rginnewpending}%
    \else
      \global\setbox\m@rginstack{#3@#4}=\vbox{\box0\unvbox\m@rginstack{#3@#4}}%
    \fi
  \fi
}

% Draw the held notes at the top of column #1, called at the start of the
% column's vbox in each output routine.
\def\m@rginstacks#1{%
  \m@rginovfout
  \ifnum0\ifvoid\m@rginstack{#1@left}\else1\fi\ifvoid\m@rginstack{#1@right}\else1\fi>0
    \vbox to 0pt{\vbox to 0pt{\unvbox\m@rginstack{#1@left}\vss}\nointerlineskip
      \vbox to 0pt{\unvbox\m@rginstack{#1@right}\vss}\vss}\nointerlineskip
  \fi
}

\newbox\m@rginboxes
\def\m@rgincollect#1{\setbox1=\vbox{#1}\trace{e}{vadj adding box \the\ht1+\the\dp1*\the\wd1}\global\setbox\m@rginboxes=\vbox{\box1\unvbox\m@rginboxes}\trace{e}{Adding vadj for \the\wd\m@rginboxes}}
\def\m@rginvadj{\ifnum\currentgrouptype=0 \ifvoid\m@rginboxes\else\trace{e}{vadjusting boxes in mode \the\currentgrouptype}\loop
    \global\setbox\m@rginboxes=\vbox{\unvbox\m@rginboxes \global\setbox0=\lastbox}%
    \ifvoid0 \tempfalse\else\temptrue\fi \iftemp\trace{e}{vadjusting a box wd=\the\wd0}\vadjust{\unvbox0}%
  \repeat
  \global\setbox\m@rginboxes=\box\voidb@x\fi\fi}

\def\pl@cemnote#1#2#3#4#5#6{%
% #1 = vbox, #2 = side, #3 = valign, #4 = distance from text, #5 = x-shift, #6 = y-shift
  \edef\t@mp{#2}\edef\t@mpv{#3}%
  \dimen9=\ht#1
  \trace{M}{\t@mpv, \t@mp \space kerns: pstrut=\the\dp\pstr@t, x-shift=\the#5, y-shift=\the#6,
            text = \the#4, side = #2, valign = #3, ht=\the\dimen9+\the\dp#1*\the\wd#1, hsize=\the\hsize}%
  \ifx\t@mp\size@SPAN\vadjust{\box #1}\else\ifm@rginheld
    % lap the note against the column, keeping its natural height and depth
    \begingroup
    \setbox0=\hbox to \hsize{\ifx\t@mp\l@ft\llap{\box #1\kern #5\kern #4}\hfil
      \else\hfil\rlap{\kern #5\kern #4\box #1}\fi}%
    \m@rginhold
    \x@\let\x@\do@djust\ifx\t@mpv\@inline\m@rgincollect\else\vadjust\fi
    \do@djust{\kern -\dimen9\kern #6\writem@rginnote\kern-#6\kern\dimen9}%
    \endgroup
  \else
    \begingroup\leftskip=0pt\rightskip=0pt\parfillskip=0pt\everypar={}%
    \dimen8=\lastdepth \dimen7=\dp#1
    \x@\let\x@\do@djust\ifx\t@mpv\@inline\m@rgincollect\else\vadjust\fi
    \setbox0=\vbox to 0pt{\let\par=\endgraf\everypar={}%
      \ifx\t@mpv\@bottom\else\ifx\t@mpv\@inline\else\vss\fi\fi
      \trace{e}{Inside (\t@mpv) \the\ht#1+\the\dp#1*\the\wd#1}%
      \ifx\t@mp\l@ft\llap{\box #1 \kern #5\kern #4}%
      \else\noindent\hfill\rlap{\kern #5\kern #4 \box #1}%
      \fi
      \ifx\t@mpv\@top\else\vss\fi}%
    \dp0=0pt
    \do@djust{\trace{e}{Inside (0) \the\ht0+\the\dp0*\the\wd0}%
      \kern -\dimen9\kern #6\writem@rginnote\box0\kern-#6\kern\dimen9
    }%
    \endgroup
  \fi\fi
}

\def\finishm@rginnotes "#1"{\immediate\closeout\m@rginnotes\m@rginnotesopenfalse
  \ifvoid\m@rginpending\else\MSG{*** marginnotes: some pushed margin notes were never drawn}\fi}

\def\openm@rginnotes "#1"{
  \ifm@rginnotesopen\else
    \trace{M}{reading #1}
    \catcode`\@=11
    \includeifpresent{#1}
    \openout\m@rginnotes="#1"
    \m@rginnotesopentrue
    \write\m@rginnotes{\percent\string\@marginnote{ref}{marker}{side}{valign}{gap}{width}{height}{depth}{xshift}{yshift}{pageno}{xpos}{ypos}}%
    \m@rginpushes
    \addtoendhooks{\finishm@rginnotes "#1"}
    \addtoeveryparhooks{\m@rginvadj}
  \fi
}

\m@rginnotesopenfalse
\def\writem@rginnote{%
  \pdfsavepos
  \x@\write\x@\m@rginnotes\x@{\wr@tet@mp{\the\pageno}{\the\pdflastxpos}{\the\pdflastypos}}\nobreak% since ther emay be baselineskip glue following
}

\newdimen\marginnotesgap     % distance from edge of text to start of notes. Not enforced here
\newdimen\marginnoteswidth   % width of marginal notes allocation - marginnotesgap
\def\m@rginnoteref#1#2{% #1 = reference, #2 = marker. Sets \mn@ref
  \edef\mn@ref{#1}%
  \x@\ifx\csname mn-lastref-#2\endcsname\mn@ref
    \x@\xdef\csname mn-refcount-#2\endcsname{\the\numexpr\csname mn-refcount-#2\endcsname+1\relax}%
    \edef\mn@ref{\mn@ref[\csname mn-refcount-#2\endcsname]}%
  \else
    \x@\global\x@\let\csname mn-lastref-#2\endcsname\mn@ref
    \x@\gdef\csname mn-refcount-#2\endcsname{0}%
  \fi
}

\def\m@rginside#1#2#3{% #1 = reference (from \m@rginnoteref), #2 = marker, #3 = style side. Sets \p@s, \p@geno, \s@de, \m@rginkey, \ifm@rginheld
  \edef\p@s{#3}%
  \ifnum\c@rrentcols=1\edef\t@mpa{page}\else\edef\t@mpa{col}\fi
  \ifcsname marginnote-#1@#2:\t@mpa no\endcsname\edef\p@geno{\csname marginnote-#1@#2:\t@mpa no\endcsname}\else
    \ifm@rginalnotesw@rn\else\MSG{*** marginnotes: no marginal note for #1-#2; re-run for accurate note positions}\global\m@rginalnotesw@rntrue\fi
    \edef\p@geno{\the\pageno}\fi
  \ifx\p@s\@uter \ifodd\p@geno \ifRTL\let\s@de\l@ft\else\let\s@de\r@ght\fi \else\ifRTL\let\s@de\r@ght\else\let\s@de\l@ft\fi\fi
  \else\ifx\p@s\@nner\ifodd\p@geno \ifRTL\let\s@de\r@ght\else\let\s@de\l@ft\fi \else\ifRTL\let\s@de\l@ft\else\let\s@de\r@ght\fi\fi
  \else\ifx\p@s\relax\let\s@de\l@ft\else\let\s@de\p@s\fi\fi\fi
  \edef\m@rginkey{#1@#2}%
  \ifcsname marginnote-\m@rginkey:push\endcsname % pushed notes take their side from the push line
    \m@rginheldtrue\def\t@mp##1##2##3{\def\s@de{##3}}\x@\x@\x@\t@mp\csname marginnote-\m@rginkey:push\endcsname
  \else\m@rginheldfalse\fi
}

% Called by \n@test@rt inside a margin note's box, before its text is set.
% Works out the reference \m@rginnoteref will give the note, without counting
% it, and the side it will be drawn on, including any push. Defines \m@rginc@l,
% which \s@tsideskips uses for inner/outer justification in place of the page
% or column parity: 1 for the right margin (left in RTL), as on an odd page.
\def\m@rginc@lset#1{% #1 = note class
  \begingroup
    \edef\mn@ref{\csname nm-ref-#1\endcsname}%
    \x@\ifx\csname mn-lastref-#1\endcsname\mn@ref
      \edef\mn@ref{\mn@ref[\the\numexpr\csname mn-refcount-#1\endcsname+1\relax]}\fi
    \m@rginside{\mn@ref}{#1}{\MarginNoteSide}%
    \xdef\m@rginc@lg{\ifx\s@de\r@ght\ifRTL0\else1\fi\else\ifRTL1\else0\fi\fi}%
  \endgroup
  \let\m@rginc@l\m@rginc@lg
}

\def\d@marginnote#1#2#3#4#5#6#7{%
% #1 = vbox, #2 = reference, #3 = marker #4 = side #5 = valign #6 = x-shift #7 = y-shift
  \m@rginnoteref{#2}{#3}%
  \m@rginside{\mn@ref}{#3}{#4}%
  \dimen0 = #6%
  \dimen3 = #7%
  \trace{M}{marginnote (#1) \mn@ref \space #3, pageno=\p@geno, side=\s@de, was #4 \ifcsname marginnote-\mn@ref @#3:colno\endcsname, col=\csname marginnote-\mn@ref @#3:colno\endcsname\fi, yshift=\the\dimen3, text offset=\the\dimen0, stystack=\styst@k}%
  %\dimen0=0pt
  \ifcsname marginnote-\mn@ref @#3:xshift\endcsname\dimen1=\csname marginnote-\mn@ref @#3:xshift\endcsname \else\dimen1=0pt \fi
  \ifcsname marginnote-\mn@ref @#3:yshift\endcsname\dimen2=\csname marginnote-\mn@ref @#3:yshift\endcsname \else\dimen2=0pt \fi
  \xdef\wr@tet@mp{\string\@marginnote{\mn@ref}{#3}{\p@s}{#5}{\the\dimen0}{\the\wd#1}{\the\ht#1}{\the\dp#1}{\the\dimen1}{\the\dimen2}}%
  \advance\dimen2 by \dimen3%
  %\showbox#1
  \pl@cemnote{#1}{\s@de}{#5}{\dimen0}{\dimen1}{\dimen2}%
}

\newcount\pcount \pcount=0
{\catcode`\~=12 \lccode`\~=32 \lowercase{
\gdef\parmkr{\setbox0=\vbox{\hbox{\ch@rstylepls{zpmkr}~\m@rker\ch@rstylepls{zpmkr}*}}%
    \advance\pcount by 1
    %\llap{\copy0}%
    \d@marginnote{0}{\the\pcount}{zpmkr}{\ifRTL right\else left\fi}{top}{0pt}{0pt}}
}
}

%\plugin@endif
